\section{Din\'amica de Perron y modo estable de memoria}

En una ruta cerrada, el cociclo cil\'indrico de Perron satisface

\begin{equation}
V_n=\varphi^n.
\label{rm:eq:cosmo-perron-V}
\end{equation}

En dimensi\'on espacial \(D=3\), la escala determinada por el cociclo es

\begin{equation}
\frac{a_n}{a_0}=V_n^{1/3}=\varphi^{n/3},
\label{rm:eq:cosmo-perron-a}
\end{equation}

mientras que el modo estable de la acci\'on cuadr\'atica es

\begin{equation}
M_n=\varphi^{-2n}=V_n^{-2}.
\label{rm:eq:cosmo-perron-M}
\end{equation}

Eliminar \(n\) entre \cref{rm:eq:cosmo-perron-a,rm:eq:cosmo-perron-M} da el
resultado central:

\begin{theorem}[Ley cil\'indrica de sexto orden para la memoria]
Para todo nivel cerrado de la realizaci\'on de Perron en \(D=3\),
\begin{equation}
\boxed{
M_n=\left(\frac{a_n}{a_0}\right)^{-6}.}
\label{rm:eq:cosmo-minus-six}
\end{equation}
El exponente \(-6\) no se ajusta: es el cociente entre el exponente estable
\(-2\) y la ra\'iz dimensional \(1/3\).
\end{theorem}

En los retornos distinguidos se obtienen

\[
\frac{a_9}{a_0}=\varphi^3,
\qquad
\frac{a_{27}}{a_0}=\varphi^9,
\qquad
\frac{a_{108}}{a_0}=\varphi^{36},
\]

y los mismos niveles verifican exactamente
\(M_n=(a_n/a_0)^{-6}\). Si se declara el par\'ametro temporal aditivo \(t_n=nt_0\), la tasa
logar\'itmica es constante:

\begin{equation}
H_\varphi
:=\frac{\Delta\log a_n}{\Delta t}
=\frac{\log\varphi}{3t_0},
\qquad
H_\varphi t_0
=0.1604039416865344824992529711\ldots.
\label{rm:eq:cosmo-Hphi}
\end{equation}

El s\'imbolo \(H_\varphi\) nombra la tasa de esta escala discreta. Su lectura
como par\'ametro de Hubble requiere el transductor cosmogr\'afico; la identidad
interna no depende de esa promoci\'on.

\section{Realizaci\'on en el sistema de Einstein--Cartan}

En un fluido de esp\'in de Einstein--Cartan, la densidad cuadr\'atica de esp\'in
obedece

\begin{equation}
s^2\propto a^{-6}.
\label{rm:eq:cosmo-EC-scaling}
\end{equation}

La \cref{rm:eq:cosmo-minus-six} suministra el exponente y el modo estable sin
consultar un rebote. Para hacer la lectura f\'isica se debe declarar una
secci\'on positiva \(s_*^2\) y un mapa

\begin{equation}
\mathcal T_{\rm EC}:M_n\longmapsto
s_n^2:=s_*^2M_n
=s_*^2\left(\frac{a_n}{a_0}\right)^{-6}.
\label{rm:eq:cosmo-EC-transducer}
\end{equation}

El mapa conserva la ley de escala, pero no determina por s\'i solo la amplitud
\(s_*^2\), la corriente de esp\'in, el signo del t\'ermino de contorsi\'on ni la
ecuaci\'on de estado. Con densidad de masa, una ecuaci\'on efectiva tiene la
forma

\begin{equation}
H^2=\frac{8\pi G}{3}
\bigl(\rho_{\rm m}-\rho_{\rm spin,m}\bigr)+\cdots,
\label{rm:eq:cosmo-EC-mass}
\end{equation}

mientras que con densidad de energ\'ia debe escribirse

\begin{equation}
H^2=\frac{8\pi G}{3c^2}
\bigl(\rho_{\rm E}-\rho_{\rm spin,E}\bigr)+\cdots.
\label{rm:eq:cosmo-EC-energy}
\end{equation}

Las dos convenciones no se mezclan. Un candidato a rebote exige, adem\'as de
la igualdad de densidades, regularidad, materia y datos de frontera.

\begin{falsifier}
La transducci\'on Einstein--Cartan falla si el funcional de escala produce un
exponente distinto de \(-6\), si \(V_n\ne\varphi^n\) en la ruta declarada, o
si el mapa f\'isico no conserva positividad y dimensiones. Elegir
\(s_*^2\) retrospectivamente para forzar un umbral de rebote no demuestra el
transductor.
\end{falsifier}

\section{Comparaci\'on estructural entre el ciclo de 108 estados y la din\'amica de Perron}

\begin{center}
\begin{tabularx}{\textwidth}{>{\bfseries}p{29mm}XX}
\toprule
& CT108--de Sitter & Perron--Einstein--Cartan\\
\midrule
Antecedente & retorno circular de orden \(108\) & cociclo cil\'indrico y modo estable\\
Selector & \(r_g=R_{108}\) & \(V_n=\varphi^n\), \(D=3\)\\
Magnitudes determinadas & radio, curvatura, temperatura y entrop\'ia & expansi\'on, tasa y diluci\'on de memoria\\
Invariante & \(G_a\hbar_a\) y radio circular & \(M_na_n^6\)\\
Estatuto & interfaz de Sitter condicionada & ley interna exacta; lectura EC condicionada\\
Criterio de refutación & radio o ecuaci\'on de fondo incompatibles & p\'erdida del exponente \(-6\) o del transductor\\
\bottomrule
\end{tabularx}
\end{center}

\begin{statusbox}
Las identidades CT108, el selector \(\theta_{108}\) y la ley
\(M_n=(a_n/a_0)^{-6}\) son \status{Teorema interno exacto}. La lectura del mismo radio
como de Sitter y la identificaci\'on de \(M_n\) con una densidad de esp\'in son
\status{Realizaciones f\'isicas condicionadas}, cada una con su propio criterio de refutación.
\end{statusbox}
